% =====================================================================
% 第 2 章  文献回顾与研究假说
% ---------------------------------------------------------------------
% 编译说明同前。
% =====================================================================

\section{文献回顾与研究假说}

与本文相关的研究可以归入三支：央行书面沟通的市场效应、文本情绪的度量方法、
货币政策的行业传导。三支文献解决的问题不同，留下的空隙也不同，本章分别梳理，
最后落到本文的定位与待检验的假说。

\subsection{央行书面沟通的市场效应}

国外文献较早把央行声明当作可度量的对象。Lucca 和 Trebbi（2009）用自动化方法解析
美联储公开市场委员会声明，考察其信息如何进入资产价格；Shapiro 和 Wilson（2022）
则从声明文本中反推政策目标的权重，把文本从“语气”推进到“偏好”。这类研究的共同
前提是，央行文字中包含着关于未来政策路径的信息，而市场会对这些信息定价。

中文文献的路径相似，但对象换成了《货币政策执行报告》。姜富伟等（2021）是其中
最直接的一篇。他们以 2001 年至 2018 年的报告为样本，用中文金融情感词典计算文本
情绪、相似度与可读性，发现情绪改善伴随发布后一日的指数正收益，相似度上升伴随波动
下降，而可读性没有显著作用。更重要的是，他们把情绪拆成“政策指引”与“宏观基本面”
两部分，发现只有前者显著，理由是报告中的宏观数据在发布前已由统计局公布，市场早已
作出反应。

沿这一思路，后续研究从不同角度补充。许魁（2023）引入分析师报告，检验其在情绪传导
中的中介作用，发现分析师情绪与报告情绪高度相关，紧缩期的传导更强。王辉等（2023）
把报告拆为前瞻性与回顾性内容，构建“市场解读偏离”指标，考察预期引导对银行系统性
风险的影响。蒋海等（2025）用机器学习方法提取前瞻性情绪，解释商业银行的流动性囤积。
王熙和郑梦圆（2021）把语调指数纳入结构向量自回归，控制信息效应后“价格之谜”消失，
从宏观侧证实了文本中的信息成分。周阳、綦岚鑫（2024）则自建适配报告的领域词典，
同时考察股票、债券与外汇市场，发现市场之间存在溢出，溢出会放大或削弱情绪效应。

这些研究给出的判断大体一致：报告文本含有信息，且信息主要来自面向未来的部分。
但结论多建立在市场总量上，行业维度往往只作为稳健性出现；样本也多在 2019 年前后
截止，2020 年之后的阶段较少被检验。

\subsection{文本情绪的度量方法}

度量方法大致有三代。第一代是词典法，把文本中的情感词按极性计数再加总。姜富伟等
（2021）使用的中文金融情感词典属于这一代，其优点是口径透明、易于复现，缺点是无法
处理语境。周阳、綦岚鑫（2024）对此作了改进，在词典中纳入否定词与程度副词，使同一
情感词在不同修饰下得到不同权重。汪禹和吴彬（2025）则从词典构建入手，用种子词加点
互信息扩充、再以形式概念分析筛选主题词，为政策文本提供了一个可复制的建库流程。

第二代是机器学习与预训练模型。纪钧尧等（2026）用 FinBERT 挖掘十二个经济体央行文本
的情绪，发现情绪对汇率波动的影响强于对汇率水平的影响；蒋海等（2025）用长短期记忆
网络标注语料，再训练分类器扩展到全样本。这一代方法在语义识别上更强，但需要标注
成本，且模型版本会影响结果。

第三代是大语言模型。廖海龙（2026）用提示词工程从报告中提取情绪因子并接入量化择时
框架；林建浩和孙乐轩（2025）系统梳理了这一路径的原理与应用，同时提醒三点局限：
输出的可复现性、对提示方式的敏感性，以及黑箱带来的可解释性问题。

本文选择词典法作为主口径，原因有二。一是与姜富伟等（2021）保持可比，便于判断差异
来自样本还是来自度量；二是词典法的误差方向较为清楚，测量误差通常导致系数向零收缩，
便于解释。为回应林建浩和孙乐轩（2025）提出的可复现性问题，本文另建一套货币政策
领域词典作稳健性，并用置换检验考察结果对随机配对的敏感性。

\subsection{货币政策的行业传导}

行业层面的研究更多集中在传统货币政策变量上。周亮（2018）用 M2、SHIBOR 与 M1-M2
划分货币周期，考察行业轮动策略的效果，发现以 M2 划分时策略表现最好。游庭玮（2026）
以美联储政策利率为解释变量，用申万一级行业的日度与周度面板估计冲击效应，行业固定
效应加行业维度聚类，并通过控制市场整体收益识别行业独立冲击，结论是高估值行业对
外部利率更敏感。王春芝（2022）用事件研究法考察货币政策公告前的股价预漂移，发现
报告发布前数个交易日已存在显著的正向异常收益。

把文本与行业结合的研究相对少。陈妍玲等（2024）用报告的情绪、相似度与可读性构建
沟通指数，考察金融业及其子行业的风险承担与风险溢出，发现前瞻性报告的文本情绪显著
提高子行业间的风险溢出。这已经是行业层面的证据，但结果变量是风险而非收益，且行业
范围限于金融业。

行业层面的一个技术困难值得单独说明。报告级情绪每期只有一个取值，若在面板中加入
季度固定效应，它与时间虚拟变量共线，系数不可识别。游庭玮（2026）的处理办法是控制
市场整体收益，以此剥离共同冲击；本文则把识别放到事件窗内，用同一时点的横截面差异
替代固定效应的作用。两种做法针对的是同一个问题：如何在不牺牲时间维度的前提下控制
共同因素。

\subsection{评述与本文的位置}

综合三支文献，可以看出一处明确的空隙。文本情绪的市场效应已被反复验证，但验证集中
在收益方向；行业维度的证据要么停留在总量，要么局限于金融业的风险指标；度量口径与
样本时期的差异又使结论难以横向比较。换言之，已有研究回答了“文本是否有信息”，
但对“信息作用于价格的哪一阶矩”“这种作用是否随时间改变”讨论不足。

本文围绕这两点展开。与姜富伟等（2021）相比，本文把样本延长至 2026 年，把被解释
变量扩展到波动，并引入制度断点；与周阳、綦岚鑫（2024）相比，本文只做股票市场，
但在分段与时期两个维度上作了更细的处理；与游庭玮（2026）相比，本文的解释变量是
文本而非政策利率，因而更接近“沟通”而非“操作”的效应。

\subsection{研究假说}

根据上述讨论，提出三个待检验的假说。

\textbf{H1（方向）}：政策指引段的文本情绪越高，报告发布后的市场收益越高。
这一假说直接对应姜富伟等（2021）的结论，用于检验本文的度量与样本能否复现既有发现。
若 H1 不成立，需要区分是度量差异、样本差异，还是效应本身不稳定。

\textbf{H2（波动）}：政策指引段的文本情绪越高，报告发布日的市场波动越低；
宏观回顾段的情绪不具有这一作用。其依据是，文本提供的是关于政策路径的信息，
信息越明确，投资者分歧越小，交易行为越趋同；而回顾段的内容在发布前已被市场知悉，
不构成新信息。

\textbf{H3（断点）}：情绪效应在 2014 年前后存在差异。2013 年的流动性紧张、2014 年
利率市场化的推进、2015 年汇率形成机制的调整，以及报告自身章节结构的改写，都可能
改变市场解读文本的方式。这一假说不预设变化的方向，只检验是否存在结构性差异。

三个假说分别对应第 5 章的三组结果：H1 对应收益方程，H2 对应波动方程与双段回归，
H3 对应交互项与断点敏感性检验。
